\documentclass[11pt]{article}
\usepackage[margin=0.9in]{geometry}
\usepackage{amsmath,amssymb,amsthm,booktabs,mathrsfs}
\usepackage[T1]{fontenc}
\usepackage{lmodern,microtype}
\usepackage[colorlinks=true,linkcolor=blue,urlcolor=blue]{hyperref}
\newtheorem{theorem}{Theorem}[section]
\newtheorem{lemma}[theorem]{Lemma}
\newtheorem{proposition}[theorem]{Proposition}
\newtheorem{corollary}[theorem]{Corollary}
\theoremstyle{remark}
\newtheorem{remark}[theorem]{Remark}
\newcommand{\norm}[1]{\left\|#1\right\|}
\begin{document}
% No preamble. Unique prefix or-. Uses \norm and standard theorem environments.
% Locality of the explicit ordinary representative, pending independent audit.

\section{Locality of the chosen ordinary-constraint representative}
\label{or-section}

Gauge equivalence by itself does not preserve locality of a chosen
potential representative. For the explicit gauge used in the ordinary
constraint transfer, however, the required locality can be proved.
Write $A_*(B)$ for the corrected fixed-slice minimizer and retain the
cumulative endpoint gauge $b_j$ and the fixed bump extension
$E_{\rm pt}$ from that transfer. Set
\begin{equation}\label{or-transfer}
 \omega(A)=\frac1i\log b_j(e^{i\eta A})^{-1},\qquad
 \lambda(A)=E_{\rm pt}\omega(A),\qquad
 T(A)=\mathscr A_\eta(A,\lambda(A)),\qquad
 A_{\rm ord}(B)=T(A_*(B)).
\end{equation}
All bases are small in the unweighted global norms. Weights are
placed on variations only. Fix $0<\alpha<1$ and one sufficiently
small positive spatial exponent $\kappa$.

\subsection{The derivative of the gauge transfer is weighted bounded}

The cumulative gauge is uniformly holomorphic in the global
potential supremum norm, has $b_j-I=O(\norm{A}_\infty)$,
and has a fixed final physical dependency radius $R_b$.
The latter follows from the endpoint product and the summable
physical reading radii of the corrected trajectory.
The same properties hold for $\omega$, with $\omega(0)=0$.
For an admissible weight $w(n)/w(m)\le e^{\kappa d(n,m)}$,
local-row Cauchy gives
\begin{equation}\label{or-omega-derivative}
 \norm{D\omega(A)V}_{\infty,w}
      \le C e^{\kappa R_b}\norm{V}_{\infty,w}.
\end{equation}
Indeed, for one coarse output row, restrict the direction to that
row's dependency neighborhood. Its unweighted supremum is bounded
by the right side divided by the output weight. No derivative is
taken of this restricted direction: $\omega$ consumes bounded
potential values. Cauchy is applied to the whole bounded row,
not coefficientwise followed by a fine-dimensional count.
The same argument proves operator-valued holomorphy in the
weighted norm on a smaller unweighted background ball.

The fixed bump interpolation satisfies
\begin{equation}\label{or-point-weighted}
 \norm{E_{\rm pt}c}_{C^{2,\alpha}_w}
          \le C_\kappa\norm{c}_{\infty,w}.
\end{equation}
Each collar reads a fixed number of nearby coarse coefficients.
Bounded smooth bump derivatives through order three, weight
comparison on that neighborhood, and sampling give this estimate,
uniformly in refinement and compatible periodic volume.
Consequently
\begin{equation}\label{or-lambda}
 \norm{D\lambda(A)V}_{C^{2,\alpha}_w}
       \le C_\kappa\norm{V}_{\infty,w},\qquad
 \norm{\lambda(A)}_{C^{2,\alpha}}
       \le C\norm{A}_\infty.
\end{equation}
The base estimate in the second inequality is unweighted.

For completeness, the fine gauge action is
\[
 \mathscr A_\eta(A,\lambda)(e)=\frac1{i\eta}\log
 (e^{i\lambda(e_-)}e^{i\eta A(e)}e^{-i\lambda(e_+)}).
\]
Writing $\ell=\lambda(e_-)$ and
$v=(\lambda(e_+)-\lambda(e_-))/\eta$, its group argument is
$e^{i\eta\operatorname{Ad}_{e^{i\ell}}A(e)}
e^{i\ell}e^{-i(\ell+\eta v)}$.
It equals identity at $\eta=0$, so its logarithm divided by
$\eta$ has a uniform analytic jet extension there.
The derivative is therefore a finite sum of products of analytic
coefficients of the base jet $(A,\lambda,D\lambda)$ and the
variation jet $(V,\mu,D\mu)$. The base coefficients have bounded
unweighted $C^{1,\alpha}$ norms. The discrete product estimate
with one weighted factor gives
\begin{equation}\label{or-gauge-derivative}
 \norm{D\mathscr A_\eta(A,\lambda)[V,\mu]}_{N_{\alpha,w}}
 \le C_\kappa\bigl(\norm{V}_{N_{\alpha,w}}
                         +\norm{\mu}_{C^{2,\alpha}_w}\bigr).
\end{equation}
A first output difference uses at most two scalar differences.
The same jet estimate controls its H\"older member; there is no
inverse-mesh loss or weighted-smallness assumption on the base.
Combining \eqref{or-lambda} and \eqref{or-gauge-derivative} proves
\begin{equation}\label{or-DT}
 \norm{DT(A)V}_{N_{\alpha,w}}
                  \le C_\kappa\norm{V}_{N_{\alpha,w}},
\end{equation}
with holomorphic dependence in this weighted operator norm.

\subsection{Responses and the new leading extension}

The corrected minimizer's weighted response theorem and the chain
rule now give
\begin{equation}\label{or-first}
 D_BA_{\rm ord}(B)=DT(A_*(B))D_BA_*(B),\qquad
 \norm{D_BA_{\rm ord}(B)b}_{N_{\alpha,w}}
                 \le C_\kappa\norm{b}_{\infty,w}.
\end{equation}
Its leading extension is
\begin{equation}\label{or-leading}
 H_{\rm ord}=DT(0)H,
\end{equation}
not H. To check the sign explicitly, let
$\beta(A)=\log b_j(e^{i\eta A})/i$ and
$B_{\rm cum}=D\beta(0)$. On the common branch $\omega=-\beta$,
and the gauge action has derivative $V-D\mu$ at zero. Hence
\begin{equation}\label{or-leading-explicit}
 DT(0)=I+DE_{\rm pt}B_{\rm cum},\qquad
 H_{\rm ord}=H+DE_{\rm pt}B_{\rm cum}H.
\end{equation}
Thus $dH_{\rm ord}=dH$.
The exact transfer identity
$\mathcal F_{\rm ord}(T(A))=\mathcal F_{\rm corr}(A)$
also gives $Q_{\rm ord}DT(0)=Q$ and
$Q_{\rm ord}H_{\rm ord}=I$, where $Q_{\rm ord}$ is the
derivative of the ordinary homogeneous constraint at zero.

The response family in \eqref{or-first} is holomorphic on a
common unweighted complex coarse ball with a uniform weighted
operator norm. Operator Cauchy on a smaller ball yields
\begin{equation}\label{or-linear-correction}
 \norm{(D_BA_{\rm ord}(B)-H_{\rm ord})b}_{N_{\alpha,w}}
       \le C_\kappa\norm{B}_\infty\norm{b}_{\infty,w}.
\end{equation}
If $\delta$ is a fixed analytic margin, then for every $n\ge1$,
\begin{equation}\label{or-higher}
 \norm{D_B^nA_{\rm ord}(B)[b_1,\ldots,b_n]}_{N_{\alpha,w}}
 \le C_\kappa(n-1)!\left(\frac e\delta\right)^{n-1}
       \norm{b_1}_{\infty,w}\prod_{i=2}^n\norm{b_i}_\infty.
\end{equation}
For $n>1$ apply multivariable Cauchy to the first-response
operator with $n-1$ unweighted directions and use
$(n-1)^{n-1}\le e^{n-1}(n-1)!$. The case $n=1$ is
\eqref{or-first}. Symmetry permits any selected bounded direction
to carry the weight.

For a direction supported in S, choose
$w_S(n)=e^{\kappa d(n,S)}$. Equations \eqref{or-first} and
\eqref{or-higher} give exponential decay away from S. If all
directions are supported in sets $S_i$, applying the bound
separately to each weighted direction gives the factor
$e^{-\kappa\max_i d(n,S_i)}$. No tree-length estimate is inferred.
Integration along a straight coarse segment gives
\begin{equation}\label{or-finite-change}
 \norm{A_{\rm ord}(B')-A_{\rm ord}(B)}_{N_{\alpha,w}}
       \le C_\kappa\norm{B'-B}_{\infty,w}.
\end{equation}
Thus finite changes supported in S have the same spatial decay.
For complete infinite fields use bounded truncated weights and
pass to the limit. For weights that decay to zero, literal
Fr\'echet responses are first read on bounded directions with
finite weighted norm; the formulas define their bounded extension
to the full weighted space. The support weights satisfy $w_S\ge1$
and require no such qualification.

This proves locality of the particular ordinary representative
defined in \eqref{or-transfer}; the field-dependent gauge has been
estimated rather than treated as harmless by gauge equivalence.
It does not assert locality of every gauge-equivalent representative,
identify the separately moving source gauge, or cover mixed
retained-history geometries. The action-minimum statement remains
the finite homogeneous small-global-field theorem proved earlier.

\end{document}
